\documentclass[12pt,a4paper]{amsart}

\usepackage{amsmath,amssymb,amsthm}
\usepackage{mathtools}
\usepackage{hyperref}
\usepackage{geometry}
\usepackage[ngerman]{babel}
\geometry{margin=2.5cm}

% -------------------------------------------------------
\newtheorem{theorem}{Satz}[section]
\newtheorem{lemma}[theorem]{Lemma}
\newtheorem{corollary}[theorem]{Korollar}
\newtheorem{proposition}[theorem]{Proposition}
\theoremstyle{definition}
\newtheorem{definition}[theorem]{Definition}
\newtheorem{remark}[theorem]{Bemerkung}

\newcommand{\Z}{\mathbb{Z}}

% -------------------------------------------------------
\begin{document}
% -------------------------------------------------------

\title{Jedes Giuga-Pseudoprim ist quadratfrei:\\ Ein kurzer Beweis}

\author{Kurt Ingwer}
\address{Specialist Mathematics Project, Version Build 44}
\email{specialist-maths@localhost}
\date{\today}

\begin{abstract}
Ein \emph{Giuga-Pseudoprim} ist eine zusammengesetzte ganze Zahl, die ein
Primalitätskriterium von Giuga~(1950) erfüllt. Wir geben einen kurzen,
in sich geschlossenen Beweis, dass jedes Giuga-Pseudoprim notwendig quadratfrei ist.
Das Argument ist elementar: Falls $p^2 \mid n$, zwingt die schwache Giuga-Bedingung
zur Konsequenz $p \mid -1$, was unmöglich ist. Dieses Ergebnis ist grundlegend für
das Studium von Giuga-Pseudoprimen und rechtfertigt, die Untersuchung von vornherein
auf quadratfreie zusammengesetzte Zahlen zu beschränken.
\end{abstract}

\maketitle

% -------------------------------------------------------
\section{Einleitung}
% -------------------------------------------------------

Im Jahre 1950 stellte Giuga \cite{Giuga1950} die Vermutung auf, dass eine natürliche
Zahl $n > 1$ genau dann prim ist, wenn $\sum_{k=1}^{n-1} k^{n-1} \equiv -1 \pmod{n}$ gilt.
Ein zusammengesetztes Gegenbeispiel würde als \emph{Giuga-Pseudoprim} bezeichnet.

Borwein, Borwein, Girgensohn und Parnes \cite{Borwein1996} charakterisierten solche
Zahlen als zusammengesetzte Zahlen, die für jeden Primteiler $p \mid n$ erfüllen:
\begin{enumerate}
  \item[\textup{(S)}] $p \mid \dfrac{n}{p} - 1$ \quad \emph{(schwache Bedingung)},
  \item[\textup{(St)}] $(p-1) \mid \dfrac{n}{p} - 1$ \quad \emph{(starke Bedingung)}.
\end{enumerate}

Die vorliegende Notiz beweist, dass jede ganze Zahl, die sogar nur die schwache
Bedingung~(S) für alle Primteiler erfüllt, quadratfrei sein muss.

% -------------------------------------------------------
\section{Hauptresultat}
% -------------------------------------------------------

\begin{theorem}
\label{thm:qfrei}
Jedes Giuga-Pseudoprim ist quadratfrei. Genauer: Wenn eine zusammengesetzte ganze Zahl~$n$
die schwache Giuga-Bedingung~\textup{(S)} für alle Primteiler $p \mid n$ erfüllt,
dann ist $n$ quadratfrei.
\end{theorem}

\begin{proof}
Angenommen, $p^2 \mid n$ für eine Primzahl~$p$. Dann gilt $p \mid \tfrac{n}{p}$, also
\[
  \frac{n}{p} \equiv 0 \pmod{p},
\]
woraus folgt $\tfrac{n}{p} - 1 \equiv -1 \pmod{p}$.
Die schwache Bedingung~(S) fordert $p \mid \tfrac{n}{p} - 1$, also $p \mid -1$.
Da jedoch $p \ge 2$, gilt $p \nmid 1$ und damit $p \nmid -1$ — Widerspruch.
Folglich kann keine Primzahl~$p$ die Bedingung $p^2 \mid n$ erfüllen, d.\,h.\ $n$ ist
quadratfrei.
\end{proof}

\begin{remark}
Der Beweis verwendet nur Bedingung~(S); die starke Bedingung~(St) wird nicht benötigt.
Damit gilt die stärkere Aussage: Jede zusammengesetzte Zahl, die nur~(S) für alle
Primfaktoren erfüllt, ist bereits quadratfrei.
\end{remark}

\begin{remark}
Die Quadratfreiheit ist eine unmittelbare Konsequenz der schwachen Bedingung:
Falls $p^2 \mid n$, ist $\tfrac{n}{p}$ durch~$p$ teilbar, also $\tfrac{n}{p} - 1$
kongruent zu~$-1$ modulo~$p$, und keine Primzahl kann~$-1$ teilen.
\end{remark}

% -------------------------------------------------------
\section{Folgerungen}
% -------------------------------------------------------

\begin{corollary}
Jedes Giuga-Pseudoprim~$n$ hat die Form $n = p_1 p_2 \cdots p_k$
mit $p_1 < p_2 < \cdots < p_k$ paarweise verschiedenen Primzahlen.
\end{corollary}

\begin{corollary}
Für ein Giuga-Pseudoprim $n = p_1 \cdots p_k$ sind die Quotienten
$\tfrac{n}{p_i}$ Produkte verschiedener Primzahlen, die von~$p_i$ verschieden sind.
\end{corollary}

Die hier bewiesene Quadratfreiheit ist der Ausgangspunkt für weitere Strukturaussagen
über Giuga-Pseudoprimen, einschließlich des Nachweises, dass keine solche Zahl weniger
als vier Primfaktoren haben kann \cite{Ingwer2026a}, sowie der zugehörigen
CRT-Widersprüche für Giuga-Carmichael-Zahlen \cite{Ingwer2026c}.

% -------------------------------------------------------
\section{Numerische Illustration}
% -------------------------------------------------------

Die bekannten \emph{Giuga-Zahlen} (die nur die schwache Bedingung erfüllen) sind
$\{30, 858, 1722, 66198, \ldots\}$. Alle sind quadratfrei:
\begin{align*}
  30 &= 2 \cdot 3 \cdot 5, \\
  858 &= 2 \cdot 3 \cdot 11 \cdot 13, \\
  1722 &= 2 \cdot 3 \cdot 7 \cdot 41, \\
  66198 &= 2 \cdot 3 \cdot 11 \cdot 17 \cdot 59.
\end{align*}
Satz~\ref{thm:qfrei} garantiert dies: Jede zusammengesetzte Zahl, die~(S) erfüllt,
ist quadratfrei. Man verifiziert direkt, dass jede dieser Zahlen die starke
Bedingung~(St) für mindestens einen Primfaktor verletzt und somit kein
Giuga-Pseudoprim ist.

% -------------------------------------------------------
\begin{thebibliography}{10}

\bibitem{Giuga1950}
G.~Giuga,
\emph{Su una presumibile propriet\`a caratteristica dei numeri primi},
Ist.\ Lombardo Sci.\ Lett.\ Rend.\ Cl.\ Sci.\ Mat.\ Nat.\ \textbf{83}
(1950), 511--528.

\bibitem{Borwein1996}
J.~M.~Borwein, P.~B.~Borwein, R.~Girgensohn und S.~Parnes,
\emph{Giuga's conjecture on primality},
Amer.\ Math.\ Monthly \textbf{103} (1996), Nr.\ 1, 40--50.

\bibitem{Ingwer2026a}
K.~Ingwer,
\emph{Keine zusammengesetzte Zahl mit genau drei Primfaktoren ist ein Giuga-Pseudoprim},
Specialist Mathematics Project (2026).

\bibitem{Ingwer2026c}
K.~Ingwer,
\emph{Über die Unverträglichkeit von Giuga-Pseudoprimen und Carmichael-Zahlen},
Specialist Mathematics Project (2026).

\end{thebibliography}

\end{document}
